\section{From a Document to a Fetch Tree}
\label{sec:ch17-from-a-document-to-a-fetch-tree}
\index{planner!how a plan is built|(}%

Planning proper is four steps: select the operation, select nodes, build
planning paths, and walk them once per datasource to produce a raw plan. What
comes out is not the plan a query-plan header shows you. Post-processing turns
it into that, and post-processing is not part of planning at all: it sits
outside the planner and the router calls it by hand immediately afterwards.
That split is not an accident of layering. Several of the things
post-processing does need the whole flat list of fetches at once, and a visitor
building that list cannot see the end of it while it is still in the middle.

\subsection*{Choosing between two subgraphs that could both answer}

\index{planner!node suggestions and the sort}%
The interesting decision is which subgraph serves a field when more than one
could. In Mosaic that happens less than you would think, because
chapter~\ref{ch:strangling-the-monolith} cut the seams so that almost every
field has one owner. It happens constantly on a graph with shareable fields and
\code{@provides}. The decision is a sort:

\begin{minted}[fontsize=\footnotesize,breaklines]{go}
	// sort by the number of jumps, selectable child count, selected siblings
	slices.SortFunc(nodesInfos, func(a, b nodeInfo) int {
		if usePriority {
			// desc on count
			if n := cmp.Compare(b.selectableChilds, a.selectableChilds); n != 0 {
				return n
			}

			if a.hasSelectedSibling != b.hasSelectedSibling {
				if a.hasSelectedSibling {
					return -1 // a comes first
				}
				if b.hasSelectedSibling {
					return 1 // b comes first
				}
			}
		}

		// acs order from 0 to n
		return cmp.Compare(a.jumpCount, b.jumpCount)
	})
\end{minted}

Three keys. First, how much of the remaining subtree this datasource could serve
if chosen, descending, so a subgraph that can answer a field and four of its
children beats one that can answer the field alone. Second, whether a sibling on
the same source has already been selected, which is the rule that keeps one
fetch from splitting into two. Third, the number of \code{@key} jumps needed to
get there, ascending.

Then it takes the first one. There is no cost model, no scoring of the finished
plan, and no comparison of alternatives: the candidates for this node are put in
an order and the head is chosen. \code{@requires} and \code{@key} both expand
through the same field-collection machinery, and \code{@provides} enters earlier,
while nodes are being collected, which is why a promised field looks to the sort
like a field the promising subgraph simply has.

\subsection*{Apollo does the opposite thing}

\index{Apollo Router!planner contrasted with Cosmo's}%
This book has mostly treated the two federation stacks as interchangeable at the
level of ideas, and here they are not. Apollo's query planner generates
candidate plans and prices them. Its own comment on the function that does it:

\begin{minted}[fontsize=\footnotesize,breaklines]{text}
 * Given some initial partial plan and a list of options for the remaining parts that need to be added to that plan to make it complete,
 * this method "efficiently" generates (or at least evaluate) all the possible complete plans and the returns the "best" one (the one
 * with the lowest cost).
\end{minted}

The cost function is explicit, and its two constants say what Apollo believes
about federated execution:

\begin{minted}[fontsize=\footnotesize,breaklines]{typescript}
const fetchCost = 1000;
\end{minted}

\begin{minted}[fontsize=\footnotesize,breaklines]{typescript}
const pipeliningCost = 100;
\end{minted}

A fetch is priced at a thousand fields, so the network hop dwarfs anything the
resolvers do; and fetches in sequence are multiplied relative to fetches in
parallel, so depth hurts far more than width. The second of those is the belief
chapter~\ref{ch:strangling-the-monolith} argued from the shape of Mosaic's plan,
when it read the \code{Parallel} node as evidence that the seams had been cut in
the right places. Here the same belief is a constant in somebody's source file.

Because the search space is a product of the choices, it needs a ceiling, and
there is one: the number of complete plans the planner will evaluate defaults to
ten thousand, and it is a configuration key with a validation error behind it.

\index{planner!optimiser against heuristic}%
So one planner is an optimiser that needs a safety valve and the other is a
heuristic that cannot run long.
Figure~\ref{fig:ch17-two-planners} puts them side by side. I do not think either
side of that trade is obviously right, and I distrust my own instinct here,
because a sort is easier to explain and being easier to explain is not an
argument. What I can say is what each one buys. Apollo will find a better plan
on a graph where the good plan is not the locally obvious one, and it will
sometimes spend real time doing it. Cosmo will produce its plan in time that
scales with the query rather than with the graph's ambiguity, and on a graph
where two routes genuinely differ in cost it has no machinery for noticing.

\begin{figure}[htbp]
  \centering
  % The two federation planners answer "which subgraph serves this field" in
% opposite styles. Referenced from chapter 17. Bare tikzpicture; the figure
% environment, caption and label live at the call site.
%
% The point of the picture is the asymmetry at the bottom of each column.
% Apollo enumerates complete candidate plans and prices each one, so a
% cartesian product needs a ceiling (maxEvaluatedPlans) to stay bounded.
% Cosmo never enumerates: at each node with more than one candidate
% subgraph it sorts by three keys and takes the first, so there is nothing
% to bound. Neither column argues the other is wrong.
\begin{tikzpicture}[
    font=\small,
    keep/.style={draw, thick, rounded corners=2pt, minimum width=26mm,
                 minimum height=8mm, align=center, font=\scriptsize,
                 fill=black!8},
    discard/.style={draw, dashed, rounded corners=2pt, minimum width=26mm,
                    minimum height=8mm, align=center, font=\scriptsize,
                    fill=black!5},
    crit/.style={draw, rounded corners=2pt, minimum width=30mm,
                 minimum height=8mm, align=center, font=\scriptsize},
    flow/.style={-{Stealth[length=2mm]}},
    note/.style={font=\scriptsize, gray, align=center, inner sep=1pt},
    hdr/.style={font=\scriptsize\itshape, align=center}
  ]

  % -- left: Apollo, enumerate and price -------------------------------------
  \node[hdr] at (2.0,4.6) {Apollo: enumerate complete\\plans, price, keep cheapest};

  \node[note] at (2.0,3.6)
    {\code{fetchCost = 1000} per fetch\\
     \code{pipeliningCost = 100} per sequential stage};

  \draw[dashed] (0.7,2.6) -- (3.3,2.6);
  \node[note, anchor=north] at (2.0,2.2)
    {ceiling: \code{maxEvaluatedPlans}\\defaults to \code{10000}};

  \node[discard] (pl1) at (2.0,1.0)  {\code{plan 1}\\higher cost};
  \node[keep]    (pl2) at (2.0,-0.4) {\code{plan 2}\\cheapest, kept};
  \node[discard] (pl3) at (2.0,-1.8) {\code{plan 3}\\higher cost};

  \draw (pl1.south west) -- (pl1.north east);
  \draw (pl3.south west) -- (pl3.north east);

  % -- right: Cosmo, sort and take the first ---------------------------------
  \node[hdr] at (9.0,4.6) {Cosmo: at each node, sort\\candidates and take the first};

  \node[note] at (9.0,3.6)
    {no cost function;\\three sort keys decide it};

  \node[note, anchor=north] at (9.0,2.2)
    {no \code{maxEvaluatedPlans}:\\nothing enumerated, nothing to bound};

  \node[crit] (s1) at (9.0,1.0)  {1. subtree served\\(descending)};
  \node[crit] (s2) at (9.0,-0.4) {2. sibling already\\selected (prefer it)};
  \node[crit] (s3) at (9.0,-1.8) {3. \code{@key} jumps\\(ascending)};
  \node[keep] (win) at (9.0,-3.2) {first in the\\sorted order};

  \draw[flow] (s1) -- (s2);
  \draw[flow] (s2) -- (s3);
  \draw[flow] (s3) -- (win);

\end{tikzpicture}

  \caption{The same job, done two ways. Apollo enumerates complete plans, prices
  each with a cost function whose constants say a fetch is worth a thousand
  fields and sequencing is worth a hundred, and keeps the cheapest, which is why
  it needs a ceiling on how many it will consider. Cosmo sorts the candidates
  for each node by three local keys and takes the first, which is why it needs
  no ceiling and has nowhere to express that one route is dearer than another.}
  \label{fig:ch17-two-planners}
\end{figure}

That difference also explains why the two projects put their caches where they
do. A planner that may evaluate ten thousand plans has a great deal to amortise.
A planner that sorts a short list per node has much less, and Cosmo caches it
anyway, which tells you the cost is not nothing.

\subsection*{What post-processing adds}

\index{planner!post-processing}%
The pass the router runs afterwards is where a raw plan becomes the thing
chapter~\ref{ch:how-a-federated-query-actually-runs} printed. It deduplicates
fetches, assigns the fetch identifiers those plans carry, resolves the input
templates, converts a generic fetch into the concrete entity or batch-entity
form, and works out the dependency ordering that decides what can run at the
same time.

That last one is worth pinning down, because \enquote{Parallel} in a printed
query plan could easily be a label somebody applies for readability. It is not.
The grouping is computed from the fetch dependency graph, and the group is run
with an ordinary Go \code{errgroup}, which is to say concurrently, for real.

\medskip
Which is a good moment to look at what those fetches actually carry, because the
router does one more thing to them on the way out that no query plan shows.

\index{planner!how a plan is built|)}%
